\section{Problem and objective}
Synthetic data offers a way to construct training examples under controlled
conditions. This project asks how to judge whether such data is suitable
for a computer vision task, rather than judging images only by their appearance.
The assessment must connect dataset properties to the intended use and make
explicit the evidence and resources needed to support a judgement.

A starting point is the framework of \textcite{alaaHowFaithfulYour2022},
which distinguishes fidelity, diversity and generalization through
$\alpha$-Precision, $\beta$-Recall and Authenticity.
This motivates examining complementary dimensions of quality.
Its applicability to particular vision tasks, and the resources required
to apply it alongside other measures, will be investigated in the review.
The project does not assume in advance that one metric or threshold is
appropriate for every task.

\section{Research question and scope}
\begin{researchbox}{Main research question}
\MainResearchQuestion
\end{researchbox}
\textbf{Subquestions}
\begin{enumerate}
\item \textbf{RQ1 --- Requirements:} \RequirementsQuestion
\item \textbf{RQ2 --- Metrics:} \MetricsQuestion
\item \textbf{RQ3 --- Assessment resources:} \ResourcesQuestion
\end{enumerate}

The review will focus on image classification, object detection and semantic
segmentation. It will consider both rendered or simulated images and learned
image generation, distinguishing synthetic-only from mixed real/synthetic
training. Domain-general evaluation methods may inform the review, but
evidence from unrelated modalities will not be treated as direct evidence
of effectiveness in vision.

\section{Research gap and expectations}
The gap to investigate is whether existing evaluation approaches provide
sufficient guidance for connecting vision-specific requirements to both
measurements and their practical prerequisites. This is a candidate gap,
not a claim that such guidance is absent. We will test its relevance against
existing frameworks and surveys and narrow it if the literature already
addresses parts of the problem.

Building on the distinction between fidelity, diversity and generalization
in Alaa et al., we expect complementary measures to be needed rather than
a single interchangeable quality score. We also expect the usefulness of
an assessment to depend on its reference data and target task. These are
provisional analytical expectations to examine for supporting and conflicting
evidence, not statistical hypotheses that this review will experimentally test.

The objective is an evidence-backed assessment framework linking each
requirement to suitable measurements, necessary inputs and limitations.
This is a literature-review project; training a new generator or conducting
a model benchmark is outside the proposed scope.
